\documentclass[11pt]{article}
\usepackage{amsmath,amssymb,amsthm}
\newtheorem{lemma}{Lemma}
\newtheorem{theorem}[lemma]{Theorem}
\newtheorem{corollary}[lemma]{Corollary}
\begin{document}
\title{A fixed line and an orbit bound}
\author{}
\date{}
\maketitle

Let $V$ be a finite-dimensional vector space over $\mathbb F_q$.
All groups below act linearly on the right. A rational line means a
one-dimensional $\mathbb F_q$-subspace. The observations below use ordinary
linear algebra and the orbit--stabilizer theorem. They make no claim of
historical priority. Their purpose here is to combine two tests in the
finite census: the exclusion of invariant lines and the exclusion of
subgroups of bounded order.

\begin{lemma}\label{lem:line}
Let $H\leq \operatorname{GL}(V)$ and let $z\in H$ have finite order $m$,
where $\gcd(m,q-1)=1$. Suppose that
$\operatorname{Fix}_V(z)=\langle v\rangle$ for a nonzero vector $v$.
If $H$ preserves a rational line, then $|v^H|\leq q-1$.
If $H$ is generated by involutions and $q$ is odd, the sharper bound
$|v^H|\leq 2$ holds.
\end{lemma}

\begin{proof}
Suppose that $W$ is an $H$-invariant rational line. The restriction of
$z$ to $W$ is multiplication by an element $\lambda\in\mathbb F_q^\times$.
Its order divides both $m$ and $q-1$, so $\lambda=1$. Consequently
$W\subseteq\operatorname{Fix}_V(z)$. Both spaces have dimension one;
therefore $W=\langle v\rangle$. Each element of $H$ sends $v$ to a
nonzero vector on this line, which contains exactly $q-1$ such vectors.
This proves the first bound.

For the second assertion, each involutory generator acts on $W$ by a
scalar whose square is one. In odd characteristic that scalar is $1$ or
$-1$. Every product of the generators therefore sends $v$ to $v$ or
$-v$. This gives the claimed bound.
\end{proof}

\begin{theorem}\label{thm:cover}
Let $G\leq\operatorname{GL}(V)$ be finite. Fix a prime $\ell$ not dividing
$q-1$ and an integer $B$. Assume the following complete cover statement:
every proper subgroup of $G$ containing an element of order $\ell$ either
preserves a rational line or is contained in a subgroup of order at most
$B$.

Let $H\leq G$ contain an element $z$ of order $\ell$, and suppose that
$\operatorname{Fix}_V(z)=\langle v\rangle$ with $v\ne0$. If $v^H$ contains
more than
\[
 \max\{q-1,\lfloor B/\ell\rfloor\}
\]
distinct vectors, then $H=G$.
\end{theorem}

\begin{proof}
Suppose that $H$ is proper. It contains $z$, so the complete cover
statement applies. In the first case, Lemma~\ref{lem:line} gives
$|v^H|\leq q-1$.

In the second case, $|H|\leq B$. The stabilizer $H_v$ contains the
actual cyclic subgroup $\langle z\rangle$, because $vz=v$. Thus
$|H_v|\geq\ell$. The orbit--stabilizer theorem gives
\[
 |v^H|=|H|/|H_v|\leq B/\ell.
\]
Since the left side is an integer, it is at most $\lfloor B/\ell\rfloor$.
Both alternatives contradict the assumed number of distinct vectors.
Hence $H$ is not proper, and $H=G$.
\end{proof}

\begin{corollary}\label{cor:g2}
Let $G$ be the acquired faithful natural matrix model of $G_2(11)$ on
$\mathbb F_{11}^{7}$. Import the complete maximal-subgroup consequence
that a proper subgroup containing $19$ lies in a rational line
stabilizer or in a subgroup isomorphic to $J_1$, of order $175560$.
Let $H\leq G$. Suppose that an actual word $z$ in the given generators
of $H$ has order $19$, that its complete fixed space has dimension one,
and that $v$ is a nonzero vector in that space. A certificate of $9241$
distinct vectors reachable from $v$ proves $H=G$.
\end{corollary}

\begin{proof}
Apply Theorem~\ref{thm:cover} with $q=11$, $\ell=19$ and $B=175560$.
We have $19\nmid10$ and $175560/19=9240$. The required threshold is
therefore $9240$, and $9241$ distinct reachable vectors exceed it.
Reachability places every certified vector in $v^H$; closure of the
supplied set is unnecessary.
\end{proof}

For an original rank-four marking, take $H$ to be either end generated
by three consecutive marks. If $H=G$, the omitted nonidentity mark belongs
to $H$. The subgroup on its singleton label then has a nontrivial
intersection with the subgroup on the three disjoint labels. Their label
sets have empty intersection. This directly violates the intersection
property. This argument applies to either end; it does not transfer a
previously proved left-end hypothesis to a right-end case.

\paragraph{Proof obligations and attribution.}
The complete cover in Corollary~\ref{cor:g2} is an attributed input from
the retained odd-characteristic maximal-subgroup classification and its
natural-module interpretation. It is the same consequence recorded in
\texttt{subgroup\_orbits/THEORY.rst}; it is not proved by a list obtained
after imposing irreducibility. The original tuple, actual membership in
$G$, and faithful matrix frame are also explicit imported inputs.

The new use in this census is the combination in
Lemma~\ref{lem:line}: a verified one-dimensional fixed space replaces a
separate search through simultaneous sign systems for an invariant line.
The computational seed verifier checks the complete kernel, rather than
only checking $vz=v$, so it supplies the additional dimension hypothesis.
A separate consumer must still verify all $9241$ distinct parent images.
A search reaching its storage limit supplies no such certificate by
itself. This note records a proof component, not a newly resolved
conjecture or a completed original manuscript.
\end{document}
